% Part: normal-modal-logic
% Chapter: frame-correspondence
% Section: introduction

\documentclass[../../../include/open-logic-section]{subfiles}

\begin{document}

\olfileid{nml}{frd}{int}

\olsection{प्रस्तावना}

मोडल तर्कशास्त्रात प्राप्यता संबंध आणि तो संबंध असलेल्या
प्रतिमानांतील काही !!{formula}s ची सत्यता यांमधील संबंध
महत्त्वाचा असतो. उदाहरणार्थ, प्राप्यता संबंध परावर्ती
आहे, म्हणजे प्रत्येक $w
\in W$ साठी $Rww$ आहे, असे समजा. दुसऱ्या शब्दांत,
प्रत्येक जग स्वतःपासून प्राप्य आहे. मग एखाद्या जगात~$w$
येथे $\Box !A$ सत्य असेल, तर ज्या प्राप्य जगांत~$!A$
सत्य असणे आवश्यक आहे त्यांत $w$ स्वतःही येते. म्हणून
$\mModel{M}$ चा प्राप्यता संबंध~$R$ परावर्ती असेल, तर
कोणतेही जग $w$ आणि कोणतेही सूत्र $!A$ घेतले, तरी
$\Box !A
\lif !A$ तेथे सत्य असेल. दुसऱ्या शब्दांत, $\Box p \lif
p$ हा रूपबंध आणि त्याचे सर्व आदेशन-दृष्टांत~$\mModel{M}$
मध्ये सत्य आहेत.

मात्र व्यत्यास असत्य आहे. उदाहरणार्थ, $\Box p \lif p$
हे~$\mModel{M}$ मध्ये सत्य असल्यावरून $R$ परावर्ती आहे
असे ठरत नाही. कारण $\Box p \lif
p$ सर्व जगांत सत्य असलेले अपरावर्ती प्रतिमान~$\mModel{M}$
सहज मिळते: स्वतःपासून प्राप्य नसलेले एकच जग~$w$ असलेले,
पण $w \in V(p)$ असलेले प्रतिमान घ्या. $p$ चे सत्यतामूल्य
योग्य निवडल्यास अपरावर्ती प्रतिमानात एखाद्या विशिष्ट
सूत्रासाठी $\Box !A \lif !A$ सत्य करता येते.

यासाठी मूल्यांकनाचा घटक~$V$ बाजूला ठेवू. $V(p)$ मध्ये
कोणती जगं आहेत याची पर्वा न करता, $\Box p \lif p$ हे
$\mModel{M}$ मधील सर्व जगांत सत्य असावे अशी अट घातल्यास
$R$ परावर्ती असणे आवश्यक ठरते. कारण अपरावर्ती
प्रतिमानात $Rww$ नसलेले किमान एक जग~$w$ असते.
$V(p) =
W \setminus \{w\}$ असे ठेवा. मग~$w$ सोडून इतर
सर्व जगांत $p$ सत्य असते; म्हणून~$w$ पासून प्राप्य
असलेल्या सर्व जगांतही ते सत्य असते ($w$ हे~$w$ पासून
प्राप्य नसल्याची खात्री आहे, आणि $p$ असत्य असलेले
$w$ हेच एकमेव जग आहे). दुसरीकडे, $w$ येथे $p$ असत्य
आहे; म्हणून $\Box
p \lif p$ हे~$w$ येथे असत्य आहे.

यावरून मूल्यांकनाशिवाय प्रतिमानाची रचना दाखवण्यासाठी
स्वतंत्र संकेतन हवे, असे सुचते; अशा रचनांना \emph{चौकटी}
म्हणतो. चौकट $\mModel{F}$ म्हणजे जगांचा संच आणि
प्राप्यता संबंध यांची जोडी $\tuple{W, R}$.
प्रत्येक प्रतिमान $\tuple{W, R, V}$ हे $\tuple{W, R}$
या चौकटीवर \emph{आधारित} असते. उलट, एखादी चौकट
तिच्यावर आधारित प्रतिमानांचा वर्ग ठरवते; आणि चौकटींचा
वर्ग त्यांपैकी कोणत्यातरी चौकटीवर आधारित प्रतिमानांचा
वर्ग ठरवतो. !!a{formula} चे चौकटीत \emph{वैध} असणे,
म्हणजे $\mModel{F} \Entails !A$, अशी व्याख्या करता येते:
चौकट~$\mModel{F}$ वर आधारित प्रत्येक $\mModel{M}$
साठी $\mSat{M}{!A}$.

या संकेतनाने !!{formula}s आणि चौकटींचे वर्ग यांतील
अनुरूपता सांगता येते: उदाहरणार्थ,
$\mModel{F} \Entails \Box p \lif p$ तेव्हा आणि केवळ
तेव्हाच, जेव्हा $\mModel{F}$ परावर्ती आहे.

\end{document}
